\subsection{Separable algebras}

DEFINITION 1. --- Let A be a commutative algebra over a field K. Then A is said to be separable over K, or also a separable K-algebra if the ring L ⊗K A is reduced for every extension L of K.

Every separable algebra is clearly reduced. For a partial converse see Th. 3 (V, p.~125).

Examples. --- 1) Let A be a polynomial algebra \(K[X_{i}]_{i \in I}\) . For every extension L of K the ring \(L \otimes_{K} A\) is isomorphic to \(L[X_{i}]_{i \in I}\) (III, p.~449, Remark 2) and hence is an integral domain (IV, p.~9, Prop. 8). In other words, every polynomial algebra over a field K is a separable K-algebra.

\begin{enumerate}
\def\labelenumi{\arabic{enumi})}
\setcounter{enumi}{1}
\item
  Let A be a commutative algebra of finite degree over a field K. For A to be separable it is necessary and sufficient for it to be etale (V, p.~34, Th. 4).
\item
  Let E be an algebraic extension of a field K. If L is an extension of K, then the ring \(L \otimes_{K} E\) is a union of the subrings \(L \otimes_{K} F\) where F ranges over the set of all subextensions of finite degree of E; therefore the ring \(L \otimes_{K} E\) is reduced if and only if this is true of \(L \otimes_{K} F\) for every subextension F of E, of finite degree over K. Taking Example 2 into account we see that E is a separable algebra in the sense of the above Def. 1 if and only if it is a separable algebraic extension in the sense of Def. 1 of V, p.~36.
\end{enumerate}

PROPOSITION 3. --- Let K be a field.

\begin{enumerate}
\def\labelenumi{\alph{enumi})}
\item
  Every subalgebra of a separable K-algebra is separable.
\item
  Every direct limit of separable K-algebras is separable.
\item
  Every product of separable K-algebras is separable.
\item
  Let A be a K-algebra and \(K'\) an extension of K. For A to be separable it is necessary and sufficient for the \(K'\) -algebra \(\mathbf{A}_{(\mathbf{K}')}\) obtained from A by extension of scalars to be separable.
\end{enumerate}

Let \(\mathbf{L}\) be an extension of \(\mathbf{K}\) . Let \(\mathbf{A}\) be a separable algebra over \(\mathbf{K}\) and \(\mathbf{B}\) a subalgebra of \(\mathbf{A}\) ; then the ring \(\mathbf{L} \otimes_{\mathbf{K}} \mathbf{A}\) is reduced, and \(\mathbf{L} \otimes_{\mathbf{K}} \mathbf{B}\) is isomorphic to a subring of \(\mathbf{L} \otimes_{\mathbf{K}} \mathbf{A}\) , hence is reduced. Thus \(\mathbf{B}\) is separable and \(a)\) is proved. In the same way \(b)\) may be proved, using the canonical isomorphism of \(L \otimes_{K} \varinjlim A_i\) with \(\varinjlim L \otimes_K A_i\) (II, p.~290, Prop. 7) and \(c)\) is proved by remarking that \(L \otimes_K \left( \prod_{i \in I} A_i \right)\) is isomorphic to a subring of \(\prod_{i \in I} (L \otimes_K A_i)\) (II, p.~306, Prop. 15).

We shall use the notation of \(d\) ). For every extension \(\mathbf{L}'\) of \(\mathbf{K}'\) the rings \(\mathbf{L}' \otimes_{\mathbf{K}'} \mathbf{A}_{(\mathbf{K}')}\) and \(\mathbf{L}' \otimes_{\mathbf{K}} \mathbf{A}\) are isomorphic (II, p.~278, Prop. 2). We conclude that if \(\mathbf{A}\) is a separable \(\mathbf{K}\) -algebra, then \(\mathbf{A}_{(\mathbf{K}')}\) is a separable \(\mathbf{K}'\) -algebra. Conversely suppose that \(\mathbf{A}_{(\mathbf{K}')}\) is a separable \(\mathbf{K}'\) -algebra; the preceding remark shows that \(\mathbf{L}' \otimes_{\mathbf{K}} \mathbf{A}\) is reduced for every extension \(\mathbf{L}'\) of \(\mathbf{K}\) containing \(\mathbf{K}'\) as subextension. Let \(\mathbf{L}\) be an extension of \(\mathbf{K}\) ; by the Scholium (V, p.~13) there exists an extension \(\mathbf{L}'\) of \(\mathbf{K}\) containing \(\mathbf{K}'\) as a subextension and a \(\mathbf{K}\) -homomorphism of \(\mathbf{L}\) into \(\mathbf{L}'\) ; the ring \(\mathbf{L} \otimes_{\mathbf{K}} \mathbf{A}\) is thus isomorphic to a subring of \(\mathbf{L}' \otimes_{\mathbf{K}} \mathbf{A}\) and hence reduced. This proves that \(\mathbf{A}\) is a separable \(\mathbf{K}\) -algebra.

PROPOSITION 4. --- Let A be a separable algebra over a field K and let B be the total ring of fractions of A; then the K-algebra B is separable.

Let S be the set of cancellable elements of A. We have identified A with a subring of B (I, p.~113); further, every element of S is invertible in B and every element of B has the form \(as^{-1}\) with \(a \in \mathbf{A}\) and \(s \in \mathbf{S}\) . Let \(L\) be an extension of K and \(x\) a nilpotent element of \(\mathbf{L} \otimes_{\mathbf{K}} \mathbf{B}\) . Then \(x\) may be written in the form \(x = \sum_{i=1}^{n} y_i \otimes a_i s_i^{-1}\) with \(y_i \in \mathbf{L}\) , \(a_i \in \mathbf{A}\) , \(s_i \in \mathbf{S}\) for \(1 \leqslant i \leqslant n\) . If we put \(s = s_1 \ldots s_n\)

then \(x(1 \otimes s)\) belongs to the subring \(L \otimes_{K} A\) of \(L \otimes_{K} B\) ; since \(x(1 \otimes s)\) is nilpotent and A separable over K, we have \(x(1 \otimes s) = 0\) , and since s is invertible in B we find that x = 0. This proves the ring \(L \otimes_{K} B\) to be reduced, whence the proposition.

PROPOSITION 5. --- Let K be a field and A, B commutative K-algebras. If A is reduced and B separable then A ⊗K B is reduced.

By the Cor. of Prop. 2 (V, p.~119) A is isomorphic to a subalgebra of a product \(\prod_{i\in I}L_i\) where \(L_{i}\) is an extension of K for \(i\in I\) . Therefore \(\mathbf{A}\otimes_{\mathbf{K}}\mathbf{B}\) is isomorphic to a

subring of \(\left(\prod_{i\in I}L_i\right)\otimes_K B\) and this latter ring is isomorphic to a subring of

\(\prod_{i\in I}(\mathbf{L}_i\otimes_{\mathbb{K}}\mathbf{B})\) (II, p.~306, Prop. 15). Since \(\mathbf{B}\) is separable, each of the rings

\(L_{i} \otimes_{K} B\) is reduced, and so the same is true of \(\prod_{i \in I} (L_{i} \otimes_{K} B)\) , and a fortiori of \(A \otimes_{K} B\) .

COROLLARY 1. --- Let K be a field, L a separable extension of K and f a polynomial in K{[}X{]}. If f has no multiple factors in K{[}X{]}, it has also no multiple factors in L{[}X{]}.

If \(f\) has no multiple factor in \(\mathbf{K}[\mathbf{X}]\) , the quotient ring \(\mathbf{K}[\mathbf{X}] / (f)\) is reduced; for \(f\) is a product of irreducible polynomials \(f_i\) relatively prime in pairs, hence

K{[}X{]}/(f) is isomorphic, by I, p.~110, Prop. 9, to the product of fields K{[}X{]}/(f\_i). By Prop. 5 the ring L{[}X{]}/(f), being isomorphic to L ⊗\_K K{[}X{]}/(f), is reduced; if g is a non-constant polynomial of L{[}X{]} such that g\^{}2 divides f, then the residue class of fg\^{}\{-1\} in L{[}X{]}/(f) is a non-zero nilpotent element; hence f has no multiple factors in L{[}X{]}.

COROLLARY 2. --- Let A and B be two commutative K-algebras. If A and B are separable, the same is true of \(\mathbf{A} \otimes_{\mathbf{K}} \mathbf{B}\) .

Let L be an extension of K. The ring \(L \otimes_{K} A\) is reduced because A is separable; now Prop. 5 shows the ring \((L \otimes_{K} A) \otimes_{K} B\) (isomorphic to \(L \otimes_{K} (A \otimes_{K} B)\) ) to be reduced, whence the Corollary.

